% paper/sections/intro.tex — Introduction section, filled in Phase 88 Plan 03 (MANU-02).
\section{Introduction}

\subsection*{Statement of need}

Functional data analysis (FDA) addresses data where the fundamental observation
is a curve, surface, or trajectory evaluated over a continuum---temperature
records over a year, spectroscopic absorption profiles, or gait acceleration
waveforms.  Rather than treating each observed value as an independent scalar,
FDA represents entire trajectories as elements of a function space, enabling
operations such as depth-based centrality ranking, elastic alignment of phase
variation, FPCA-based dimension reduction, and scalar- or function-on-function
regression.
Practitioners working in Python have relatively few mature options.
\texttt{scikit-fda}~0.10.1 \citep{ramoscarreno_scikitfda_2024} is the strongest
existing Python peer: it provides representation and basis smoothing, elastic
registration, functional depth and outlier detection, dense FPCA, clustering,
classification, and scalar-on-function, function-on-scalar, concurrent and
historical function-on-function regression.
However, as documented in Table~\ref{tab:comparison}, it does not cover
functional time series forecasting, statistical process monitoring, conformal
prediction and tolerance bands, density or Fr\'{e}chet analysis of
object-valued data, or sparse PACE FPCA; it offers no grounded parameter advisor and no
scientific-provenance layer.
\texttt{FDApy}~1.0.3 \citep{golovkine_fdapy_2025} focuses on representation,
multivariate FPCA, and simulation, and does not address the remaining families
at all.
In R, the relevant methods are distributed across multiple packages: \texttt{fda}
\citep{r_fda,ramsay_hooker_graves_2009} covers representation, smoothing,
registration, and dense FPCA; \texttt{fda.usc} \citep{febrerobande_fdausc_2012}
adds depth, classification, scalar-on-function regression, and hypothesis
tests; \texttt{refund} \citep{r_refund} provides scalar- and
function-on-function regression and several FPCA variants;
\texttt{funData} \citep{happkurz_fundata_2020} and \texttt{tidyfun}
\citep{r_tidyfun}, with its companion \texttt{tf} \citep{r_tf}, offer
data-structure infrastructure, and \texttt{tf} adds depth and registration.
No single R package spans all the method families listed in
Table~\ref{tab:comparison}, and coordinating across these packages requires
manual data conversion at each interface boundary.
The Matlab toolkits fdaM \citep{ramsay_hooker_graves_2009} (last release 2017)
and PACE \citep{yao_muller_wang_2005} (v2.17, 2015; its methods have since been
ported in part to the R package \texttt{fdapace}) are no longer actively
developed and cannot be used in Python workflows without external bridges.

Several individual areas are, however, served well by specialised packages:
\texttt{fdapace} \citep{r_fdapace} for sparse FPCA, \texttt{ftsa}
\citep{r_ftsa} for functional time series, \texttt{funcharts}
\citep{capezza_funcharts_2023} for functional control charts,
\texttt{frechet} \citep{r_frechet} for Fr\'{e}chet regression,
\texttt{fdasrvf} \citep{r_fdasrvf} for elastic analysis, and \texttt{roahd}
and \texttt{fdaoutlier} \citep{r_roahd,r_fdaoutlier} for depth and outlier
detection.

\texttt{fdars} does not aim to supersede these tools in their own niches.
Its contribution is integration: to our knowledge it is the only Python
library that covers all fourteen capability areas of
Table~\ref{tab:comparison} behind one data model and one API---\npubliccallables{}
public callables in \nsubmodules{} submodules---with scikit-learn
compatibility, a grounded parameter advisor, and a scientific-provenance map.
We make no performance claims; the design section (\S\ref{sec:design})
describes the architecture qualitatively, and \S\ref{sec:validation} checks
numerical agreement with independent implementations, including several of
the packages above.

\subsection*{Background: functional data analysis}

Functional data analysis treats data as realisations of a random process
$X(t)$, $t \in \mathcal{T}$, where each observation is an element of an
infinite-dimensional function space rather than a finite-dimensional vector
\citep{ramsay_silverman_2005,ramsay_dalzell_1991}.
In practice, observations are recorded at discrete evaluation points and
represented via basis expansions (B-splines, Fourier), kernel smoothing, or
directly as dense grids; when observations are sparse and irregular, methods
such as PACE \citep{yao_muller_wang_2005} recover the latent smooth trajectory
from noisy longitudinal measurements.
Functional principal component analysis (FPCA) reduces dimension by projecting
onto a data-adaptive orthonormal basis of eigenfunctions, playing the role that
PCA plays in multivariate analysis \citep{ramsay_silverman_2005}.
Elastic registration \citep{srivastava_klassen_2016,marron_ramsay_sangalli_srivastava_2015}
separates amplitude from phase variation by aligning curves under a group
of diffeomorphic reparametrisations before downstream inference.
Further families---including regression, depth, clustering, functional time
series \citep{hormann_kokoszka_2010,hyndman_ullah_2007}, and process
monitoring---extend classical multivariate techniques to the functional
setting; \citet{ferraty_vieu_2006} and \citet{ramsay_silverman_2005} survey
the field comprehensively.

\subsection*{Contributions and paper outline}

The rest of the paper covers: the software design and Rust/PyO3 architecture
(\S\ref{sec:design}); the data-representation model (\S\ref{sec:represent});
a capability tour by method family, with executed, script-sourced code snippets
(\S\ref{sec:capabilities}); a comparison with related software
(\S\ref{sec:comparison}); the grounded advisor and scientific-provenance map
(\S\ref{sec:advisor}); four case studies on real data
(\S\ref{sec:casestudies}); validation, including agreement with independent
implementations (\S\ref{sec:validation}); limitations
(\S\ref{sec:limitations}); and availability (\S\ref{sec:availability}).
